\documentclass[11pt]{article}
\usepackage{coursenotes}
\begin{document}
\notesheader{5}{Evaluating targeting policies}{What is the list worth, measured on data it was not chosen on?}

\begin{decision}
A model produces a ranked list of customers to target. Before it is deployed, someone must say what the list is
worth, in money, with a standard error, and whether it beats the list already in use. That number has to
come from randomized data the list was not chosen on.
\end{decision}

\begin{goals}
  \item define a policy's value and gain, and estimate them from a randomized sample by inverse-probability
  weighting with the design's own probabilities;
  \item read gain, profit, and Qini curves, and summarize a ranking with RATE;
  \item separate ranking from calibration, say which one a decision uses, and recalibrate a score by isotonic
  regression without changing its ranking;
  \item explain why a policy's value on the data that chose it is optimistic, and run the freeze-and-evaluate
  protocol;
  \item compare two policies with a paired standard error.
\end{goals}

%=====================================================================================
\sectionone

\subsection{Policies, values, and gains}

A \emph{policy} is a rule $\pi: x \mapsto \{0,1\}$. Let $R(d)$ be the reward a unit would generate under decision $d$, a
potential outcome in money. The policy's \emph{value} and its \emph{gain} over treating no one are
\[
  V(\pi) = \E[R(\pi(X))], \qquad G(\pi) = V(\pi) - V(\mathbf{0}) .
\]
With margin $m$ per purchase, cost $c$ per redemption, and $k$ per unit treated, $R(d) = (m - cd)Y(d) - kd$, and
the net value of treating $x$ is $v(x) = \E[R(1) - R(0) \mid X = x] = m\tau(x) - c\mu_1(x) - k$.

\begin{prop}{The policy-value identity}{pvi}
$G(\pi) = \E[\pi(X)\,v(X)] = \Pr(\pi(X) = 1)\cdot\E[v(X) \mid \pi(X) = 1]$.
\end{prop}
\begin{proof}
$R(\pi(X)) - R(0) = \pi(X)\{R(1) - R(0)\}$; take expectations, conditioning on $X$ first.
\end{proof}

The gain depends only on the net values of the units the policy selects, not on how well the score predicts
outcomes. It is reach times average effect: a policy can raise its average by treating fewer units and still lose
total value. Since $v(x)$ is never observed for anyone, a list's gain is a population average and must be
estimated as one.

\subsection{Estimating value from a randomized sample}

\begin{prop}{Inverse-probability-weighted value}{ipw}
Suppose $D \indep \{R(0), R(1)\} \mid X$ with known $e(x) \in (0,1)$, and let $p_\pi(X) = e(X)$ if $\pi(X) = 1$ and
$1 - e(X)$ otherwise. Then
\[
  V(\pi) = \E\!\left[\frac{\ind{D = \pi(X)}}{p_\pi(X)}\,R\right], \qquad
  G(\pi) = \E\!\left[\pi(X)\left\{\frac{DR}{e(X)} - \frac{(1-D)R}{1-e(X)}\right\}\right].
\]
\end{prop}
\begin{proof}
On $\{D = \pi(X)\}$, $R = R(\pi(X))$. Given $X$, independence gives
$\E[\ind{D = \pi(X)}R(\pi(X)) \mid X] = p_\pi(X)\E[R(\pi(X)) \mid X]$; divide and average over $X$. The gain follows the
same way from $\E[DR/e(X) \mid X] = \E[R(1) \mid X]$ and $\E[(1-D)R/(1-e(X)) \mid X] = \E[R(0) \mid X]$.
\end{proof}

The sample version averages \emph{row contributions}
$g_i = \pi(X_i)\{D_iR_i/e(X_i) - (1-D_i)R_i/(1-e(X_i))\}$, with standard error $\text{sd}(g_i)/\sqrt n$. With $e$ known,
the identity is exact and the only uncertainty is sampling. The weights must be the design's: in a 25\%-treated
test a treated buyer counts $1/0.25 = 4$ times, not twice. Normalizing the weights within each arm (the
\emph{H\'ajek} version) is usually less variable.

\begin{prop}{Variance of a row contribution}{var}
$\Var(g_i) = \E\!\left[\pi(X)\left\{\dfrac{\E[R(1)^2 \mid X]}{e(X)} + \dfrac{\E[R(0)^2 \mid X]}{1-e(X)}\right\}\right] - G(\pi)^2$.
\end{prop}
\begin{proof}
$D(1-D) = 0$ removes the cross term, so $g_i^2 = \pi(X_i)\{D_iR_i^2/e^2 + (1-D_i)R_i^2/(1-e)^2\}$; condition on $X$ and
use independence, then subtract $G(\pi)^2$.
\end{proof}

The variance is driven by the second moments of the whole reward, not by the effect: a small effect on a large,
noisy reward needs a large evaluation sample. Replacing $g_i$ with the doubly robust score (subtract a prediction
of the reward fitted on other data, then weight the residual) keeps the expectation and can cut the variance
sharply.

\subsection{Gain, profit, and Qini curves}

Rank units by a score $s(x)$ and let $S_\phi$ be the top fraction $\phi$. The \emph{gain curve}
$Q(\phi) = \phi\{\E[Y(1) \mid S_\phi] - \E[Y(0) \mid S_\phi]\}$ is the incremental outcome per unit of the population from
treating the top $\phi$. Every ranking starts at $0$ and ends at the ATE; random ranking traces $\phi\cdot\ATE$. With
four equal groups whose effects are $0.20, 0.10, 0.00, -0.10$ (ATE $0.05$), the best ranking gives $Q = 0.050, 0.075, 0.075,
0.050$ at $\phi = \tfrac14, \tfrac12, \tfrac34, 1$: it rises, flattens, and falls once it reaches the harmed group. A curve
that does not start at $0$ and end at the ATE has been computed incorrectly. The \emph{Qini curve} is the same plot in counts.

\begin{prop}{Estimate rates within a group, then scale}{rates}
For a group $S$ of $n_S$ units with $n_{dS}$ in arm $d$ and arm totals $B_{dS}$, the estimator
$n_S(B_{1S}/n_{1S} - B_{0S}/n_{0S})$ is unbiased, given arm sizes, for the group's incremental outcome. The raw count
difference $B_{1S} - B_{0S}$ is not, unless the arms are equal, and even then it is half of it.
\end{prop}
\begin{proof}
Each arm is a random sample of the group, so $B_{dS}/n_{dS}$ is unbiased for its mean; scale by $n_S$. The count
difference has expectation $n_{1S}\mu_{1S} - n_{0S}\mu_{0S}$.
\end{proof}

Replace the outcome by net value and the curve becomes the \emph{profit curve} $P(\phi) = \phi\,\E[v(X) \mid S_\phi]$.

\begin{prop}{Where the profit curve peaks}{peak}
If units are ranked by their true net value $v$ (with a continuous distribution), $P$ is concave, its slope at $\phi$
is the net value of the marginal unit, and it is maximized where that net value is zero.
\end{prop}
\begin{proof}
$P(\phi) = \int_0^\phi v_{(u)}\,du$ with $v_{(u)}$ the net value at rank $u$, which decreases in $u$; its derivative
$v_{(\phi)}$ is decreasing and crosses zero at the maximum.
\end{proof}

A gain curve in purchases can still be rising while the profit curve falls: the marginal unit adds purchases but
costs more than they earn. Under a per-redemption cost, ranking by $\hatt$ is not ranking by $v$.

\paragraph{RATE.} The \emph{targeting operator characteristic} $\text{TOC}(q) = \E[\tau(X) \mid s(X)$ in the top $q] - \ATE$
measures how much better the top $q$ respond than average. The \emph{rank-weighted average treatment effect} is
\[
  \text{RATE} = \int_0^1 w(q)\,\text{TOC}(q)\,dq .
\]
With $w(q) = 1$ it is the AUTOC, which emphasizes the top of the list; with $w(q) = q$ it is the Qini coefficient,
which weighs the whole list. Under a ranking independent of the effects, TOC is zero at every $q$ and so is any
RATE; a RATE significantly above zero on held-out data shows the score ranks better than random. Use AUTOC when
only a small fraction will be treated and the Qini coefficient when a large fraction will. In practice the integral
is approximated on a grid (for example, deciles).

\subsection{Ranking is not calibration}

\begin{prop}{Gain curves see only ranks}{ranks}
For any strictly increasing $h$, the scores $s$ and $h \circ s$ give identical gain curves and RATEs. A threshold
policy $\ind{s(x) > t}$ and its gain change with $h$.
\end{prop}
\begin{proof}
The groups $S_\phi$ depend only on the order of scores, which $h$ preserves; which units exceed a fixed $t$ does
change.
\end{proof}

Under a fixed budget (``treat the top $B$'') only the ranking matters. Under a break-even threshold
(``treat if $\hatt(x) > k/m$'') the score must be calibrated: a model whose scores are twice the true effects ranks
perfectly and sends to everyone with true effect above half the break-even. Check calibration with the test and
group effects of Meeting~4.

\paragraph{Recalibrating a score.} A miscalibrated score can be repaired without touching its ranking. On units not used
to fit the score, let $\psi_i = D_iY_i/e(X_i) - (1 - D_i)Y_i/\{1 - e(X_i)\}$, or the doubly robust score $\phi_i$ of Meeting~3; either
has conditional mean $\tau(X_i)$.

\begin{prop}{Isotonic calibration}{iso}
Let $s$ be a score fixed before the calibration sample, and let $\hat h$ minimize $\sum_i\{\psi_i - h(s(X_i))\}^2$ over
non-decreasing functions $h$. Then $\hat h \circ s$ never reverses the order of $s$, and $\hat h$ is constant on blocks of
consecutive scores, equal on each block to the block's average of $\psi_i$.
\end{prop}
\begin{proof}
A non-decreasing $h$ cannot reverse an order. Take the blocks to be maximal runs of units with a common fitted value, so
neighboring blocks differ strictly. If a block's value differed from its average $\psi$, moving it slightly toward the
average would keep the order with its neighbors and lower the sum of squares.
\end{proof}

Each block's value estimates that block's average effect, a group effect on data-chosen groups, so the cautions of
Meeting~4 apply. By Result~\ref{prop:ranks}, gain curves and RATE are unchanged except for the order within tied blocks;
what changes is the scale a threshold rule reads. The pool-adjacent-violators algorithm computes $\hat h$: start from each
unit (or score bin) alone and merge any neighbor pair whose averages are out of order until none is. Fit the map on the
selection sample of the protocol below, never on the evaluation sample.

\subsection{Choosing and evaluating are different jobs}

\begin{prop}{The chosen policy's selection value is optimistic}{optimistic}
Let $\hat G_S(\pi_1), \dots, \hat G_S(\pi_K)$ be unbiased on a selection sample $S$ and $\hat\pi$ the maximizer. Then
$\E[\hat G_S(\hat\pi)] \ge \E[G(\hat\pi)]$. On an evaluation sample $E$ independent of everything used to choose $\hat\pi$,
$\E[\hat G_E(\hat\pi) \mid \hat\pi] = G(\hat\pi)$.
\end{prop}
\begin{proof}
The first is the winner's curse (Meeting~4). For the second, $\hat\pi$ is fixed given data independent of $E$, so
Result~\ref{prop:ipw} applies to it on $E$.
\end{proof}

Candidates are not only models: cut-offs, feature sets, budgets, and any adjustment after looking at results are
all choices. An automated system that tunes its own cut-off is making one.

\paragraph{The evaluation protocol.} (1) \emph{Split before fitting}: draw training, selection, and evaluation samples
from the experiment at random. (2) \emph{Fit} on training data. (3) \emph{Choose} among candidates on the selection
sample. (4) \emph{Freeze}: record the chosen policy's decision for every evaluation unit, and the code, before
evaluation outcomes are released. (5) \emph{Evaluate once} and report that number. A disappointing evaluation is a
correct result, not a reason to return to step 3, which would turn the evaluation sample into a selection sample.

\subsection{Comparing two policies}

\begin{prop}{Paired comparison}{paired}
For frozen policies $\pi_a, \pi_b$ evaluated on the same randomized sample, let $d_i = g_i(\pi_a) - g_i(\pi_b)$. Then
$\bar d$ is unbiased for $G(\pi_a) - G(\pi_b)$ with standard error $\text{sd}(d_i)/\sqrt n$.
\end{prop}
\begin{proof}
Linearity and Result~\ref{prop:ipw} for each policy; the $d_i$ are functions of independent rows.
\end{proof}

$d_i = 0$ wherever the two lists agree. Combining the two lists' separate SEs as if they were independent ignores
the shared noise and, for overlapping lists, badly overstates the uncertainty of the difference.

\begin{pitfall}[evaluation errors]
Reporting a policy's value on the data that chose it; reusing $e = 0.5$ for a design that was not 50/50; drawing
gain curves with different endpoints; comparing two lists with unpaired SEs; re-thresholding after seeing the
evaluation result.
\end{pitfall}

%=====================================================================================
\sectiontwo

\paragraph{Setting.} QuickBite emails customers an offer to join QuickBite Plus (a delivery-fee subscription), with a
small incentive. Each email and incentive costs ¥1.00; each additional subscription is worth ¥120 of margin.
Break-even lift $= 1/120 = 0.0083$. An uplift model was trained and its cut-off chosen on last year's experiment. This
year a 20{,}000-customer randomized test (50/50) is the evaluation sample.

\paragraph{Step 1: by hand on ten rows.} Policy: treat the four highest scores.
\begin{center}
\begin{tabular}{ccccccc}
\toprule
Customer & Score & $\pi$ & $D$ & $Y$ & $D = \pi$? & $\pi \cdot (2D-1)\cdot 2Y$ \\
\midrule
1 & 0.09 & 1 & 1 & 1 & yes & 2 \\
2 & 0.08 & 1 & 0 & 0 & no & 0 \\
3 & 0.07 & 1 & 1 & 0 & yes & 0 \\
4 & 0.06 & 1 & 0 & 0 & no & 0 \\
5 & 0.04 & 0 & 1 & 1 & no & 0 \\
6 & 0.03 & 0 & 0 & 1 & yes & 0 \\
7 & 0.02 & 0 & 1 & 0 & no & 0 \\
8 & 0.01 & 0 & 0 & 1 & yes & 0 \\
9 & 0.00 & 0 & 1 & 0 & no & 0 \\
10 & $-0.01$ & 0 & 0 & 0 & yes & 0 \\
\bottomrule
\end{tabular}
\end{center}
Agreeing rows have outcomes $1, 0, 1, 1, 0$, so $\hat V(\pi) = 3/(10 \times 0.5) = 0.6$; the five controls give
$\hat V(\mathbf 0) = 2/(10 \times 0.5) = 0.4$; gain $0.2$, equal to the average of the last column, $2/10$
(Result~\ref{prop:ipw} in outcome units).

\paragraph{Step 2: deciles on the full evaluation sample.} Deciles of 2{,}000 (1{,}000 per arm), effects by difference
in means, scaled by Result~\ref{prop:rates}:
\begin{center}
\begin{tabular}{cccccc}
\toprule
Deciles & Effect of & Incremental & Net & Mean effect & TOC \\
treated & this decile & sales & profit (¥) & of treated & \\
\midrule
1 & 0.050 & 100 & 10{,}000 & 0.0500 & 0.0365 \\
2 & 0.035 & 170 & 16{,}400 & 0.0425 & 0.0290 \\
3 & 0.025 & 220 & 20{,}400 & 0.0367 & 0.0232 \\
4 & 0.015 & 250 & 22{,}000 & 0.0313 & 0.0178 \\
5 & 0.010 & 270 & 22{,}400 & 0.0270 & 0.0135 \\
6 & 0.005 & 280 & 21{,}600 & 0.0233 & 0.0098 \\
7 & 0.002 & 284 & 20{,}080 & 0.0203 & 0.0068 \\
8 & 0.000 & 284 & 18{,}080 & 0.0178 & 0.0043 \\
9 & $-0.002$ & 280 & 15{,}600 & 0.0156 & 0.0021 \\
10 & $-0.005$ & 270 & 12{,}400 & 0.0135 & 0.0000 \\
\bottomrule
\end{tabular}
\end{center}
Net profit is $120 \times$ incremental sales $- 1.00 \times 2{,}000 \times$ deciles treated. Emailing everyone nets ¥12{,}400;
profit peaks at five deciles (¥22{,}400), where the marginal decile's effect ($0.010$) is the last above $0.0083$
(Result~\ref{prop:peak}).

\begin{center}
\begin{tikzpicture}
\begin{axis}[width=0.7\linewidth, height=5cm, xlabel={share of customers emailed}, ylabel={net profit (¥ thousand)},
  xmin=0, xmax=1, ymin=0, ymax=25, axis lines=left, legend style={draw=none, font=\footnotesize, at={(0.98,0.05)}, anchor=south east}]
\addplot[thick, sbgreen, mark=*, mark size=1.2pt] coordinates {(0,0) (0.1,10) (0.2,16.4) (0.3,20.4) (0.4,22) (0.5,22.4) (0.6,21.6) (0.7,20.08) (0.8,18.08) (0.9,15.6) (1,12.4)};
\addplot[thick, dashed, extgrey] coordinates {(0,0) (1,12.4)};
\legend{uplift model, random order}
\end{axis}
\end{tikzpicture}

\small Figure 1. The profit curve peaks at five deciles; random order earns the straight line.
\end{center}

\paragraph{Step 3: RATE.} Averaging the TOC column over the ten grid points gives AUTOC $\approx 0.0143$; weighting each by
$q$ gives a Qini-type RATE of $0.0046$. Both are positive: the score ranks better than random.

\paragraph{Step 4: selection versus evaluation.} On last year's selection data the top half showed 340 incremental
subscriptions per 20{,}000; on the evaluation sample, 270. The 70-subscription gap (about ¥8{,}400) is the optimism of
Result~\ref{prop:optimistic}. Report ¥22{,}400.

\paragraph{Step 5: calibration.} Had QuickBite used ``email if score $> 0.0083$'' with scores twice the true effects, it
would email everyone with true effect above $0.0042$ (deciles 1--6) and earn ¥21{,}600, ¥800 less. Isotonic calibration on
the selection sample (Result~\ref{prop:iso}) would have mapped each decile's score to an estimate of its own effect, about
half the raw score, and restored the five-decile rule without refitting the model.

\begin{exercise}[computation]
QuickBite compares ``top 5 deciles'' with ``top 4 deciles'' on the same evaluation sample; the two lists differ only on
decile 5, where 4\% of emailed and 3\% of non-emailed customers subscribed. With reward $R(d) = 120\,Y(d) - 1.00\,d$ and
$e = 0.5$: (a) list the possible row contributions $d_i = g_i(\pi_a) - g_i(\pi_b)$ on decile-5 rows and their mean;
(b) compute the paired SE of the per-customer difference over all $n = 20{,}000$ rows; (c) is the fifth decile worth
adding?
\end{exercise}
\answer{(a) Emailed subscriber $2(120 - 1) = 238$; emailed non-subscriber $-2$; non-emailed subscriber $-240$; non-emailed
non-subscriber $0$. Mean on decile-5 rows: $0.5(0.04 \times 238 - 0.96 \times 2) - 0.5(0.03 \times 240) = 0.20$, which is
$120 \times 0.01 - 1$. (b) $\E[d_i^2]$ on decile-5 rows $= 0.5(0.04 \times 238^2 + 0.96 \times 4) + 0.5(0.03 \times 240^2) \approx 1{,}999$;
over all rows $d_i = 0$ on 90\% of them, so $\text{sd}(d) \approx \sqrt{0.1 \times 1{,}999} = 14.1$ and
$\SE = 14.1/\sqrt{20{,}000} = 0.10$ per customer. (c) The difference is $0.02$ per customer (¥400 in total), a fifth of
one SE: not established. The reward's noise, not the effect, sets the SE (Result~\ref{prop:var}). Adding the decile is a
judgment call; a larger test would be needed to settle it.}

\begin{exercise}[computation]
On the selection sample, six equal-sized score bins (lowest first) have mean scores $0.004, 0.008, 0.012, 0.016, 0.020, 0.024$ and
mean transformed outcomes $\bar\psi = 0.002, 0.006, 0.004, 0.010, 0.018, 0.016$. (a) Run pool-adjacent-violators to get the
calibrated value of each bin. (b) With break-even $0.0083$, which bins does ``email if score $> 0.0083$'' treat, before and after
calibration? (c) Why fit the map on the selection sample rather than the evaluation sample?
\end{exercise}
\answer{(a) Bins 2 and 3 are out of order ($0.006 > 0.004$): pool to $0.005$. Bins 5 and 6 likewise: pool to $0.017$. Result:
$0.002, 0.005, 0.005, 0.010, 0.017, 0.017$, now non-decreasing. (b) Raw scores treat bins 3--6; calibrated values treat bins 4--6.
Bin 3's raw score $0.012$ overstates an effect of about $0.005$, below break-even. The order of the bins is unchanged. (c) The map
is fitted to outcomes, so it is a choice; fitting it on the evaluation sample would make that sample a selection sample
(Result~\ref{prop:optimistic}).}

\begin{exercise}[judgment]
After the evaluation, the team notices that moving the cut-off from five to four and a half deciles would raise the
estimated profit by ¥300 and proposes to report the higher figure. What should they report, and what would justify
using the new cut-off?
\end{exercise}
\answer{Report ¥22{,}400 for the frozen five-decile policy. Choosing the cut-off on the evaluation data makes it a selection
sample (Result~\ref{prop:optimistic}); the ¥300 is partly noise. To adopt the new cut-off, freeze it and evaluate it on
fresh data, such as next month's randomized holdout.}

%=====================================================================================
\sectionthree

\begin{extension}[Doubly robust policy evaluation]
Replace $R$ in the row contributions with the AIPW score: a prediction $\hat\mu_d(X)$ fitted on other data plus the weighted
residual. The estimate stays unbiased with known $e$ and its variance falls by as much as the reward model explains.
\end{extension}

\begin{extension}[Inference for RATE and for learned policies]
Yadlowsky et al.\ derive the sampling distribution of RATE estimators and justify a half-sample bootstrap. Imai and Li
(2023) give variance estimators for individualized rules evaluated on experimental data, including rules learned by
cross-fitting on the same experiment.
\end{extension}

\subsection*{Annotated reading}
\begin{readinglist}
\reading{van der Laan, Ulloa-P\'erez, Carone, and Luedtke (2023), ``Causal isotonic calibration for heterogeneous treatment
effects'', \emph{ICML}}{Isotonic calibration of effect scores on doubly robust outcomes}{Sections 1--3}{3}
\reading{Yadlowsky, Fleming, Shah, Brunskill, and Wager, ``Evaluating treatment prioritization rules via rank-weighted
average treatment effects'', \emph{JASA} (2025)}{TOC, AUTOC, and Qini as members of one family, with inference}{Sections
1--3}{3}
\reading{Imai and Li (2023), ``Experimental evaluation of individualized treatment rules'', \emph{JASA}}{Finite-sample
evaluation of a rule on a randomized sample}{Sections 1--3}{3}
\reading{Radcliffe (2007), ``Using control groups to target on predicted lift'', \emph{Direct Marketing Analytics
Journal}}{The original Qini curve, from practice}{All}{1}
\reading{Athey and Wager (2021), ``Policy learning with observational data'', \emph{Econometrica}}{Doubly robust value
estimation as the basis for choosing policies (Meeting~6)}{Sections 1--2}{3}
\end{readinglist}

\end{document}
